\chapter{Divisors and Principal Divisors}
\label{ch:divisors}

A divisor records zeros and poles as a formal integer combination of points.

\begin{definition}
For a compact Riemann surface \(X\), a divisor is a finite sum
\[
  D=\sum n_p[p],\qquad n_p\in\Z.
\]
Its degree is \(\degdiv D=\sum n_p\).  If \(f\) is meromorphic, its principal
divisor is
\[
  \divisor(f)=\sum_p \ord_p(f)[p].
\]
\end{definition}

Principal divisors are the geometric version of logarithmic derivatives.  Since
\(\Res_p(df/f)=\ord_p(f)\), the compact residue theorem gives:

\begin{proposition}
\label{prop:principal-degree-zero}
Every principal divisor on a compact Riemann surface has degree zero.
\end{proposition}

\begin{proof}
The meromorphic one-form \(df/f\) has residue \(\ord_p(f)\) at \(p\).  The sum of
residues on a compact Riemann surface is zero, so
\(\degdiv\divisor(f)=0\).
\end{proof}

\begin{definition}
The degree-zero divisor group is \(\Div^0(X)\).  The principal divisors form a
subgroup \(\Prin(X)\subset\Div^0(X)\), and
\[
  \Pic^0(X)=\Div^0(X)/\Prin(X).
\]
\end{definition}

The group \(\Pic^0(E)\) is not decorative notation.  It is the geometric home of
the same group already seen analytically as \(\C/\Lambda\).  A point \([P]-[O]\)
is a degree-zero divisor class, and the map \(P\mapsto[P]-[O]\) will identify the
curve with its Jacobian.

\section*{Exercises}
\addcontentsline{toc}{section}{Exercises}

\exerciseintent{Practice divisor bookkeeping, including the contribution at infinity.}
\begin{exercise}
Compute the divisor of \(z^2(z-1)^{-3}\) on \(\CP^1\).  First compute the
orders at \(0\) and \(1\).  Then use the coordinate \(w=1/z\) to compute the
order at \(\infty\).
\end{exercise}

\exerciseintent{Verify that principal divisors form a degree-zero subgroup.}
\begin{exercise}
For meromorphic functions \(f,g\), compute \(\divisor(fg)\),
\(\divisor(f^{-1})\), and \(\divisor(c)\) for \(c\in\C^*\).  Use these
calculations, together with degree zero for principal divisors, to show that
\(\Prin(X)\) is a subgroup of \(\Div^0(X)\).
\end{exercise}

\exerciseintent{Check the basic degree-zero divisor attached to a pointed curve.}
\begin{exercise}
Compute the degree of \([P]-[O]\) by adding its two coefficients.
\end{exercise}

\exerciseintent{Compute addition of degree-zero divisor classes in \(\Pic^0\).}
\begin{exercise}
Show that
\[
  ([P]-[O])+([Q]-[O])=[P]+[Q]-2[O]
\]
defines an element of \(\Pic^0(X)\) by explicitly computing the degree of the
right-hand side and then passing modulo \(\Prin(X)\).
\end{exercise}

\exerciseintent{Translate equality in \(\Pic^0\) into a principal-divisor condition.}
\begin{exercise}
Let \(D,D'\in\Div^0(X)\).  Compute the difference \(D-D'\) in the quotient
\(\Div^0(X)/\Prin(X)\).  State and prove the criterion for \(D\) and \(D'\) to
define the same element of \(\Pic^0(X)\).
\end{exercise}
